\section{Conclusiones}
% Versión original por si el profesor prefiere:
\begin{comment}
La depuración de la base de datos mejoró los modelos entrenados, aumentando su rendimiento en tareas de seguimiento y superando ligeramente la versión desplegada por WildSense en condiciones de buena visibilidad. Incluso cuando no se igualó totalmente a la versión anterior, se alcanzaron métricas competitivas con un menor número de instancias seguidas, lo que indica mayor selectividad del modelo y potencial reducción de la carga computacional en uso real.

La búsqueda de hiperparámetros contribuyó a optimizar el entrenamiento y permitió que modelos como Deepfish alcanzaran su máximo rendimiento, llegando a superar el SOTA reportado en trabajos previos.

La inclusión de imágenes de fondo aporta información útil pero en escasa medida, además, puede ser contraproducente al aumentar el coste de entrenamiento por mejoras marginales. Además, es crucial definir con precisión el conjunto de validación para evaluar si las imágenes de fondo realmente mejoran el desempeño, como se observó en el caso de Salmons.

Complementariamente, el \textit{Transfer Learning} puede ser una alternativa efectiva ante limitaciones de hardware o tiempo, aunque suele implicar una ligera pérdida de rendimiento. En estas tareas, SGD tiende a ofrecer mayor estabilidad y robustez frente a AdamW. Respecto a la exportación con TensorRT, FP16 es un compromiso equilibrado que acelera la inferencia sin sacrificar demasiado rendimiento, mientras que INT8 solo es recomendable en escenarios con restricciones computacionales severas o si la velocidad de inferencia es una prioridad.

\subsection{Trabajo futuro}
Queda por explorar si realmente los nuevos modelos son más eficientes en producción, midiendo su rendimiento en un entorno real y evaluando la calidad de las estimaciones de masa y volumen realizadas con sus predicciones. Se recomienda explorar otros algoritmos de seguimiento además de DeepSORT y ByteTrack, cosa que podría dar luces sobre hasta qué punto las detecciones de los modelos entrenados permiten un mejor \textit{tracking}. También es prioritario definir de forma rigurosa qué se considera una "instancia de interés" para salmones y documentar un protocolo de etiquetado que minimice ambigüedades y facilite la reproducibilidad al momento de entrenar y evaluar los modelos.
\end{comment}

\begin{itemize}
  \setlength{\itemsep}{0pt}%
  \setlength{\parskip}{0pt}%
  \setlength{\parsep}{0pt}%
    \item Se depuró programáticamente la base de datos Salmons, mejorando el seguimiento de salmones de interés en el contexto de jaulas de cultivo y la calidad de las segmentaciones en modelos entrenados. Incluso cuando no se igualó a la versión anterior, se alcanzaron métricas competitivas con un menor número de instancias seguidas, lo que indica mayor selectividad del modelo y potencial reducción de la carga computacional en uso real.
    \item La búsqueda de hiperparámetros permitió que modelos entrenados con Deepfish alcanzaran o incluso superaran el SOTA previo, lo que resalta la efectividad de este enfoque.
    \item Incluir imágenes de fondo puede aportar información, pero también degradar el rendimiento y aumentar el coste; su utilidad depende de un conjunto de validación bien definido.
  % TODO lenguaje: “costo” o “coste” unificar; preferir uno en toda la memoria.
    \item Para esta tarea SGD mostró mayor estabilidad y robustez que AdamW.
    \item El uso de \textit{Transfer Learning} incurre solo en una ligera pérdida de rendimiento, pero puede suponer una alternativa útil bajo restricciones de hardware o tiempo.
    \item Para despliegue, TensorRT-FP16 ofrece buen equilibrio precisión--latencia; INT8 solo cuando priman la velocidad o las restricciones computacionales.
\end{itemize}

\subsection{Trabajo futuro}
Como áreas de interés para trabajo futuro, se sugieren las siguientes líneas de acción:
\begin{itemize}
  \setlength{\itemsep}{0pt}%
  \setlength{\parskip}{0pt}%
  \setlength{\parsep}{0pt}%
    \item Validar en producción, midiendo impacto en \textit{throughput} y contrastar calidad de estimación de biomasa/volumen con modelos empleados actualmente.
  % TODO lenguaje: Traducir “throughput” por “tasa de procesamiento” o “rendimiento”.
    \item Evaluar \textit{trackers} adicionales más allá de DeepSORT y ByteTrack.
    \item Definir con rigor qué se considera una "instancia de interés" para salmones y documentar un protocolo de etiquetado que minimice ambigüedades y facilite la reproducibilidad al momento de entrenar y evaluar los modelos.
    \item Explorar la técnica de optimización \textit{knowledge distillation} para mejorar el rendimiento de modelos YOLO usando arquitecturas más precisas como base.
    \item Explorar el uso de arquitecturas de modelos de Segmentación de Instancias en Video (VIS), que resuelven simultáneamente segmentación y seguimiento.
\end{itemize}